\documentclass[UTF8,12pt,a4paper]{ctexart}
\usepackage[margin=2.4cm]{geometry}
\usepackage{amsmath,amssymb,graphicx,booktabs,listings,xcolor,hyperref,fancyhdr,setspace,float}
\hypersetup{hidelinks}
\setstretch{1.18}
\setlength{\parindent}{2em}
\setlength{\parskip}{3pt}
\pagestyle{fancy}
\fancyhf{}
\fancyhead[L]{图像处理课程作业}
\fancyhead[R]{运动模糊 PSF}
\fancyfoot[C]{\thepage}
\setlength{\headheight}{15pt}
\renewcommand{\headrulewidth}{0.4pt}
\lstdefinestyle{matlab}{language=Matlab,basicstyle=\ttfamily\footnotesize,keywordstyle=\bfseries,
  commentstyle=\itshape\color{gray!80!black},numbers=left,numberstyle=\tiny\color{gray},
  stepnumber=1,numbersep=7pt,frame=single,rulecolor=\color{gray!45},breaklines=true,
  showstringspaces=false,tabsize=4,columns=fullflexible,keepspaces=true}

\title{运动模糊点扩散函数的离散实现}
\author{姓名：蔡斌\quad 班级：数学jd2401\\[4pt]
学号：1024001338}
\date{2026 年 11 月 2 日}

\begin{document}
\maketitle

\section{运动模型}
课上推导了匀速直线运动的连续核。本次考虑运动方向，并比较离散核与 MATLAB 的
\lstinline|fspecial('motion',LEN,THETA)|。LEN 表示运动长度，单位为像素；THETA 从直角坐标系的水平正方向起算，逆时针为正。因此，$0^\circ$ 是水平运动，$90^\circ$ 是竖直运动。

把运动中点放在原点，单位方向向量为 $(\cos\theta,\sin\theta)$。连续线段为
\begin{equation}
  \mathbf r(s)=(s\cos\theta,s\sin\theta),\qquad -\frac{L}{2}\le s\le \frac{L}{2}.
\end{equation}
理想 PSF 沿这条线段均匀分布：
\begin{equation}
 h(x,y)=\frac{1}{L}\int_{-L/2}^{L/2}
 \delta(x-s\cos\theta)\,\delta(y-s\sin\theta)\,\mathrm ds.
\end{equation}
连续模型确定了运动轨迹，但没有规定像素网格上的边界采样和端点权重。这些离散细节会改变最终矩阵。

\section{离散核的构造}
对一个像素中心 $(x,y)$，沿运动方向的坐标为 $q=x\cos\theta+y\sin\theta$，到运动直线的垂直距离为 $d_{\perp}=|-x\sin\theta+y\cos\theta|$。到有限线段的距离为
\begin{equation}
 d=\sqrt{d_{\perp}^{2}+\max\left(|q|-\frac{L}{2},0\right)^2}.
\end{equation}
可以按距离给像素分配权重，再归一化。但 MATLAB 还要确定核尺寸、端点处理和镜像顺序。R2023b 的实现先生成最小半核，再计算距离；端点附近改用欧氏距离，随后转换成一像素宽的线性权重、镜像并归一化。代码保留了其中的 \texttt{eps} 修正项。\footnote{MathWorks, \emph{fspecial: Create predefined 2-D filter}, \url{https://www.mathworks.com/help/images/ref/fspecial.html}.}

\section{MATLAB 对照}
AI 初版按到有限线段的距离给方阵赋权。对 LEN=9、THETA=45$^\circ$，初版支持域为 $11\times11$，MATLAB 输出为 $7\times7$。差异来自核尺寸的确定方式：MATLAB 根据角度生成最小半核网格。

最终实现依次计算半核边界、端点欧氏距离和像素权重，再镜像展开并归一化。水平、竖直运动分别输出行向量和列向量。

\subsection*{AI 交互}
交互从题目中的“直接编写运动 PSF”开始。AI 给出距离采样初版，MATLAB 对照显示 $45^\circ$ 核尺寸不同；随后围绕半核网格和端点处理改写函数，再用多组 LEN、THETA 复测。

\subsection*{对话截图}
\begin{center}
\fbox{\parbox[c][3.0cm][c]{0.90\textwidth}{\centering 截图 1：原始提问与 AI 初版}}
\vspace{3mm}

\fbox{\parbox[c][3.0cm][c]{0.90\textwidth}{\centering 截图 2：尺寸差异与修改分析}}
\vspace{3mm}

\fbox{\parbox[c][3.0cm][c]{0.90\textwidth}{\centering 截图 3：最终函数与 MATLAB 测试结果}}
\end{center}

\clearpage
\section{函数实现}
\lstinputlisting[style=matlab]{motion_psf.m}

\clearpage
\section{结果验证}
测试程序将函数与本机 MATLAB R2023b 的 \texttt{fspecial} 比较。每组先检查矩阵尺寸，再用 \texttt{isequal} 检查元素。测试参数为：
\begin{itemize}
  \item LEN：$[0.25,1,1.25,2.5,5,9,13.7,20]$；
  \item THETA：$[-180,-90,-45,0,30,45,89.5,90,135,180,225,270,360]$ 度。
\end{itemize}
共比较 $8\times13=104$ 组，另检查无参数调用和省略 THETA 的默认调用。

\begin{center}
\begin{tabular}{lcc}
\toprule
检查内容 & 通过数量 & 最大绝对误差 \\
\midrule
LEN/THETA 参数组合 & 104/104 & 0 \\
默认参数调用 & 2/2 & 0 \\
\bottomrule
\end{tabular}
\end{center}

LEN=9 时几个方向的输出如下。系数和略小于 1，是 MATLAB 归一化时加入浮点修正项所致；本函数与内置输出逐元素相同。
\begin{center}
\begin{tabular}{ccc}
\toprule
THETA & 矩阵大小 & 系数和 \\
\midrule
$0^\circ$ & $1\times9$ & 0.99999999999999822 \\
$45^\circ$ & $7\times7$ & 0.99999999999999833 \\
$90^\circ$ & $9\times1$ & 0.99999999999999833 \\
\bottomrule
\end{tabular}
\end{center}

图中把核居中放入 $9\times9$ 矩阵显示；测试比较的是未补零的原始输出。
\begin{figure}[H]
  \centering
  \includegraphics[width=0.78\textwidth]{截图/psf_kernels.png}
  \caption{LEN=9 时 THETA 为 $0^\circ$、$45^\circ$ 和 $90^\circ$ 的离散核}
\end{figure}

验证环境：MATLAB R2023b（23.2.0.2365128）。\par
\noindent 测试脚本：\texttt{test\_motion\_psf.m}。

\end{document}
